% ============================================================
% Methodology appendix/section -- Corrected 1D spectral validation
% Thesis: Frequency Helmholtz Transfer Operator
% Author: F.E. Kiewiet de Jonge
% ============================================================
%
% Suggested packages if this file is used standalone:
%   \usepackage{amsmath, amssymb}
%   \usepackage{booktabs}
%   \usepackage{graphicx}
%   \usepackage{xcolor}
%
% This file is written as a thesis-ready methodology section. It can be
% included as a subsection of the numerical methodology chapter or adapted
% into an appendix.

\section{One-dimensional spectral validation methodology}
\label{sec:methodology_1d_spectral}

The two-dimensional Helmholtz systems studied in this thesis are too large for
a complete eigenvalue decomposition: a $512 \times 512$ grid contains
$262{,}144$ unknowns before accounting for complex arithmetic. A direct
eigendecomposition of the full matrix is therefore neither computationally
practical nor necessary for validating the spectral mechanism of the proposed
warm start. The one-dimensional problem provides a controlled intermediate
setting. It preserves the essential ingredients of the numerical method,
namely an indefinite Helmholtz operator, absorbing boundary layers, a shifted
Laplacian preconditioner, and a frequency-transfer neural operator, while
remaining small enough that all eigenvalues and eigenvectors can be computed
exactly.

The goal of the one-dimensional study is threefold. First, it verifies that
the eigenvalue computation itself agrees with the analytical Dirichlet
spectrum. Second, it separates the role of the Perfectly Matched Layer (PML)
from the physical interior modes. Third, it tests whether the learned
frequency-transfer operator reduces the components of the initial error that
are associated with difficult Helmholtz modes, and whether this reduction
translates into fewer Krylov iterations.


% ============================================================
\subsection{Grid, physical interior, and notation}
\label{subsec:methodology_1d_grid}
% ============================================================

The one-dimensional grid contains
\[
    N = 512
\]
unknowns. The outer
\[
    n_{\mathrm{pml}} = 112
\]
points on each side form the absorbing layer, leaving
\[
    N_{\mathrm{int}} = N - 2 n_{\mathrm{pml}} = 288
\]
physical interior unknowns. The frequency set is
\[
    \omega \in \{16,32,64,128\}.
\]
For the Dirichlet spectral analysis the grid spacing is chosen as
\[
    h = \frac{1}{N+1},
\]
corresponding to $N$ unknowns between two fixed Dirichlet endpoints outside
the computational grid. This convention is important because it gives a
closed-form reference spectrum against which the numerical eigensolver can be
validated.

The central object is the high-frequency linear system
\begin{equation}
    A_H u_H = f,
    \label{eq:oned_high_frequency_system}
\end{equation}
where $A_H = A(\omega_H)$ is the discretized Helmholtz operator at the target
frequency. A learned transfer operator $T_{\mathrm{up}}$ maps a low-frequency
solution $u_L$ to an initial guess
\begin{equation}
    x_0 = T_{\mathrm{up}}(u_L)
    \label{eq:oned_warm_start}
\end{equation}
for the high-frequency Krylov solve.


% ============================================================
\subsection{Dirichlet reference operator}
\label{subsec:methodology_1d_dirichlet}
% ============================================================

Eigenvalue component weighting is performed only for one-dimensional
Dirichlet problems. This restriction is deliberate. The Dirichlet operator is
real symmetric, so its eigenvectors form an orthonormal basis. Consequently
modal coefficients have a direct interpretation as Euclidean projections,
without the conditioning difficulties that arise for non-normal PML matrices.

The Dirichlet Helmholtz operator is written with the sign convention
\begin{equation}
    A_D(\omega) = -\frac{d^2}{dx^2} - \omega^2.
    \label{eq:oned_dirichlet_continuous}
\end{equation}
The second-order finite-difference stencil is
\begin{equation}
    (A_D u)_i =
    -\frac{u_{i-1}-2u_i+u_{i+1}}{h^2} - \omega^2 u_i,
    \label{eq:oned_dirichlet_stencil}
\end{equation}
or equivalently
\begin{equation}
    (A_D)_{i,i} = \frac{2}{h^2} - \omega^2,
    \qquad
    (A_D)_{i,i\pm1} = -\frac{1}{h^2}.
    \label{eq:oned_dirichlet_matrix_entries}
\end{equation}
For $n$ Dirichlet unknowns, the analytical eigenvalues are
\begin{equation}
    \lambda_k
    =
    \frac{4}{h^2}
    \sin^2\!\left(\frac{\pi k}{2(n+1)}\right)
    - \omega^2,
    \qquad k=1,\ldots,n.
    \label{eq:oned_dirichlet_analytic_eigs}
\end{equation}
This formula predicts a small low-index negative region followed by a
predominantly positive spectrum. Reproducing this distribution numerically is
the first validation step. If the numerical eigenvalues disagree with
\eqref{eq:oned_dirichlet_analytic_eigs}, then any subsequent spectral
interpretation is unreliable.

Two Dirichlet bases are used:
\begin{table}[h]
\centering
\begin{tabular}{lll}
\toprule
Basis & Size & Purpose \\
\midrule
$V_{288}$ & $288$ & Primary physical-interior component weighting \\
$V_{512}$ & $512$ & Full-grid Dirichlet comparison without PML eigenvectors \\
\bottomrule
\end{tabular}
\caption{Dirichlet eigenvector bases used for one-dimensional component
weighting. The $288$-point basis is the default because it represents the
physical interior. The $512$-point basis is included as a full-grid comparison
that still avoids non-normal PML eigenvectors.}
\label{tab:oned_dirichlet_bases}
\end{table}

Both bases are computed with a Hermitian eigensolver. The eigenvectors are
explicitly normalized so that
\begin{equation}
    \|v_k\|_2 = 1
    \label{eq:oned_unit_eigenvectors}
\end{equation}
for every mode. This condition is checked numerically and reported with the
component-weighting results.


% ============================================================
\subsection{Flux-form PML operator for data generation and GMRES}
\label{subsec:methodology_1d_flux_pml}
% ============================================================

The PML is not used as the modal basis for component weighting, but it is used
in the finite-difference solves and in the GMRES evaluation. This distinction
is essential. The Krylov method must solve the same full $512$-point problem
used by the numerical simulator, including the absorbing boundary. However,
modal coefficient claims should be made in a clean Dirichlet basis.

The PML is introduced by the complex stretch
\begin{equation}
    s(x) = 1 + i\frac{\sigma(x)}{\omega}.
    \label{eq:oned_pml_stretch}
\end{equation}
The corrected one-dimensional PML operator is discretized in flux form:
\begin{equation}
    A_{\mathrm{PML}} u
    =
    -\frac{1}{s}
    \frac{d}{dx}
    \left(
        \frac{1}{s}\frac{du}{dx}
    \right)
    - \omega^2 u.
    \label{eq:oned_flux_pml_operator}
\end{equation}
This form is preferred over a row-scaled stencil because it more faithfully
represents the stretched-coordinate differential operator. The PML damping
profile is
\begin{equation}
    \sigma_i
    =
    \sigma_0(\omega)
    \left(
        \frac{n_{\mathrm{pml}} - i}{n_{\mathrm{pml}}}
    \right)^p,
    \qquad p=2,
    \label{eq:oned_pml_profile}
\end{equation}
inside each absorbing layer and zero in the physical interior. The default
values are
\begin{table}[h]
\centering
\begin{tabular}{cc}
\toprule
$\omega$ & $\sigma_0(\omega)$ \\
\midrule
16  & 42.5 \\
32  & 85.0 \\
64  & 120.0 \\
128 & 180.0 \\
\bottomrule
\end{tabular}
\caption{Default PML damping strengths used in the corrected one-dimensional
flux-PML experiments. Sensitivity analysis varies these values by a scalar
factor.}
\label{tab:oned_sigma_values}
\end{table}

The resulting matrix is complex and non-Hermitian. Its eigenvalues are useful
as a diagnostic of boundary absorption and non-normality, but its right
eigenvectors are not used for coefficient weighting. This avoids attributing
physical meaning to boundary-localized PML modes.


% ============================================================
\subsection{Training variants for the frequency-transfer operator}
\label{subsec:methodology_1d_training}
% ============================================================

The learned operator $T_{\mathrm{up}}$ is trained to map a low-frequency
solution to a high-frequency initial guess,
\[
    T_{\mathrm{up}} : u_{\omega_L} \mapsto u_{\omega_H},
    \qquad \omega_H > \omega_L.
\]
The training data consist of random interior Gaussian sources. For each source
field $f$, the corrected flux-PML systems are solved at $\omega_L$ and
$\omega_H$. The solution is normalized by the root-mean-square magnitude of
the low-frequency solution on the physical interior. This prevents PML
amplitudes from dominating the scaling.

Two training variants are prioritized:
\begin{table}[h]
\centering
\begin{tabular}{llll}
\toprule
Variant & Data & Loss region & Inference rule \\
\midrule
\texttt{flux\_int} &
Corrected FD/PML &
288-point physical interior &
Zero PML strip \\
\texttt{flux\_full} &
Corrected FD/PML &
Full 512 grid &
Keep full prediction \\
\bottomrule
\end{tabular}
\caption{Primary one-dimensional training variants. The \texttt{flux\_int}
variant is the main methodology because it learns from PML-consistent data
while optimizing the physical interior. The \texttt{flux\_full} variant tests
whether learning the absorbing layer improves raw iteration counts.}
\label{tab:oned_training_variants}
\end{table}

The main thesis interpretation is based on \texttt{flux\_int}. This is the
cleanest variant: the data are generated by the same full PML solver used in
evaluation, but the objective function only rewards accuracy in the physical
domain. The PML entries of the warm start are zeroed before GMRES. This avoids
injecting spurious boundary-layer structure while still allowing the full PML
operator to absorb outgoing waves during the solve.

The \texttt{flux\_full} variant is included as a performance comparator. If it
achieves fewer GMRES iterations, this indicates that learning boundary-layer
structure can be computationally useful. However, its modal interpretation is
less clean because part of the network capacity is spent matching the PML
solution rather than only the physical wave field.


% ============================================================
\subsection{CSL-preconditioned GMRES evaluation}
\label{subsec:methodology_1d_csl}
% ============================================================

All warm starts are evaluated by solving the high-frequency PML system
\eqref{eq:oned_high_frequency_system} with GMRES or flexible GMRES. The
preconditioner is the complex shifted Laplacian (CSL)
\begin{equation}
    M_{\mathrm{CSL}}
    =
    A_H - i\beta \omega_H^2 I,
    \label{eq:oned_csl_preconditioner}
\end{equation}
with default shift parameter $\beta = 0.3$. Each Krylov iteration applies
$M_{\mathrm{CSL}}^{-1}$ to the residual. The preconditioner therefore changes
the spectrum of the operator seen by GMRES,
\begin{equation}
    M_{\mathrm{CSL}}^{-1} A_H,
    \label{eq:oned_preconditioned_operator}
\end{equation}
but it does not change the true eigenvalues of $A_H$ itself.

The warm start and CSL preconditioner affect different parts of the algorithm.
The warm start changes the initial error
\[
    e_0 = u_H - x_0,
\]
whereas the CSL preconditioner changes how GMRES propagates and reduces that
error. The central hypothesis is that a useful $T_{\mathrm{up}}$ reduces the
components of $e_0$ associated with difficult eigenmodes, so that the
preconditioned Krylov iteration reaches the tolerance in fewer steps.

The primary evaluation metric is the mean iteration count required to reach
a relative residual tolerance of $10^{-6}$. Residual curves are reported
together with the initial relative error, because a low residual and a low
solution error need not coincide at iteration zero for indefinite Helmholtz
systems.


% ============================================================
\subsection{Eigenvalue component weighting}
\label{subsec:methodology_1d_component_weighting}
% ============================================================

For a normalized Dirichlet eigenbasis $V = [v_1,\ldots,v_n]$, the initial
error is decomposed as
\begin{equation}
    e_0 = \sum_{k=1}^{n} c_k v_k,
    \qquad
    c_k = \langle v_k, e_0 \rangle.
    \label{eq:oned_component_weighting}
\end{equation}
Because $\|v_k\|_2=1$ and the basis is orthonormal, $|c_k|$ measures the
amplitude of the error in mode $k$. The cold start has $x_0=0$, so
\[
    e_0^{\mathrm{cold}} = u_H.
\]
For a warm start,
\[
    e_0^{\mathrm{warm}} = u_H - T_{\mathrm{up}}(u_L).
\]
The modal plots compare $|c_k^{\mathrm{cold}}|$ and
$|c_k^{\mathrm{warm}}|$ as a function of ordered eigenmode index $k$.

This analysis is performed for both the $288$-point physical-interior
Dirichlet basis and the $512$-point full-grid Dirichlet basis. The $288$-point
basis is the primary result because it corresponds to the physical domain on
which the loss is evaluated. The $512$-point basis is a secondary comparison
that tests whether the same qualitative modal reduction is visible on the full
grid without introducing PML eigenvectors.

The expected evidence for a successful warm start is not uniform reduction of
all modes. Instead, the important question is whether the warm start reduces
the error in the modes that are difficult for GMRES. These are typically modes
near the origin of the Helmholtz spectrum, where the operator is close to
singular and the Krylov method is most sensitive to phase and amplitude errors.


% ============================================================
\subsection{Sensitivity analysis plan}
\label{subsec:methodology_1d_sensitivity}
% ============================================================

After the baseline corrected experiment is established, sensitivity analysis
varies the numerical choices that can affect both the PML quality and the
iteration count:
\begin{table}[h]
\centering
\begin{tabular}{lll}
\toprule
Parameter & Default & Purpose of sweep \\
\midrule
$\sigma$ scale & $1.0$ & Test PML damping strength \\
PML power $p$ & $2$ & Test damping ramp shape \\
CSL shift $\beta$ & $0.3$ & Test preconditioner strength \\
Frequency pair & $16\to32$ first & Test transfer difficulty \\
Component basis & $288$ first & Compare physical and full-grid modal views \\
\bottomrule
\end{tabular}
\caption{Sensitivity parameters for the corrected one-dimensional pipeline.}
\label{tab:oned_sensitivity_parameters}
\end{table}

The sensitivity study is staged deliberately. First, the $16\to32$ transfer
case is used to verify that the corrected PML data, training loss, component
weighting, and GMRES evaluation are mutually consistent. Only after this
baseline is understood are the higher frequency pairs and broader PML/CSL
sweeps launched. This prevents an expensive sweep from obscuring a basic
methodological error.


% ============================================================
\subsection{Summary of methodological safeguards}
\label{subsec:methodology_1d_safeguards}
% ============================================================

The one-dimensional methodology contains four safeguards against misleading
spectral conclusions.
\begin{enumerate}
    \item The Dirichlet eigensolver is validated against the analytical formula
    \eqref{eq:oned_dirichlet_analytic_eigs}.
    \item Eigenvalue component weighting is performed only in unit-norm
    Dirichlet eigenvector bases.
    \item The full PML operator is used for solves and GMRES, but not for
    modal coefficient claims.
    \item Iteration-count improvements are interpreted together with initial
    solution error and modal error coefficients.
\end{enumerate}
Together, these safeguards make the one-dimensional study a reliable bridge
between analytical spectral expectations and the large two-dimensional
experiments where complete eigendecomposition is impossible.

